\section{Money and Trade}

\begin{frame}{Quantitative Theory of Money}
\showmytikz[1,2]{../shared/tikz/tikz-quantitytheoryofmoney}
\mycomment{
    Imagine a two-sector economy, with only fishermen and coconut farmers.

    The (exchange) value of fish in coconut will depend entirely on how many fish and how many coconut are brought to the market (per market day). This in turn will depend on preferences of the consumers, and on the effort needed for fishing and harvesting. 

    Imagine a seasonal shortage of coconut relative to plenty of fish. Then some fishermen will change profession and start harvesting coconut. 
    
    If labour is mobile (and all that counts), then costs to fish and harvest should equal out.

    This is a world of exchange without money, which may work if there are only fish and coconut. 
    
    As commerce expands beyond fish and coconut, a common denominator will become desirable. 
}
\end{frame}

\begin{frame}{Quantitative Theory of Money}
\showmytikz[1,2,3,4]{../shared/tikz/tikz-quantitytheoryofmoney}
\mycomment{
        Now our islanders agree using some sea shells (of an extinct species) as medium of exchange. 
        
        Farmers buy fish with shells, and fishermen with shells buy coconut. 
    
        The value of the shells (money) depends only on two factors:
        \begin{itemize}
            \item the speed of circulation
            \item the quantity of shells in circulation
        \end{itemize}
}
\end{frame}

\begin{frame}{Quantitative Theory of Money}
\showmytikz[1,3]{../shared/tikz/tikz-quantitytheoryofmoney}
\mycomment{
        For the system (or better: this model!) to work, the shells must retain circulation. For they do not have any use value, fishermen would not accept shells without farmers accepting shells.
        
        So why shells? And where do these shells come from? One quality the shells have is that nobody can create them at will. Really? But we can create debt, can we not?  
}
\end{frame}

\begin{frame}{Quantitative Theory of Money}
\showmytikz[1,3,5]{../shared/tikz/tikz-quantitytheoryofmoney}
\mycomment{
        In our model world, neither $A$ nor $B$ ever invest in new fishing nets or ladders. If a farmer wants to take time off to make a new ladder instead of harvesting, they would need to get some credit to not starve while they do this. The fishermen would have to supply shells as a loan (and then fish in exchange for shells). Where would that credit come from, if the farmers need all their money for transactions?

        In other words: the quantitative theory of money, while absolutely worth knowing and considering, has its limits. It assumes there is a world without credit, without debt, without interest rate, without past or future, without inequality, without power and ownership, without permissions, competing currencies, etc.
}
\end{frame}


\begin{frame}{The Four Functions of Money}
    \mytext{
    
    The four functions of money are: 
    \begin{enumerate}
        \item medium of exchange
        \item unit of account
        \item store of value
        \item standard of deferred payment
    \end{enumerate}
 
    }
    \mycomment{
    The state is the institution that issues currency through the central bank. As legal tender, citizens are required to accept it unit for unit (singleness of money) to settle any debt - and in particular, tax obligations.

    Money is more than just functions (1) and (2), which may be explained by the quantitative theory of money. But its function as credit (and debit) for deferred payment, as loans for investment, as deposits, as liabilities in accounts, and as store of value threatened by inflation are studied in \textbf{macroeconomics.}
    }
\end{frame}

\begin{frame}{Why Trade?}
    \mytext{
    \begin{itemize}
        \item Nobody has to buy, and
        \item nobody has to sell.
    \end{itemize}

    I could grow my own food, make my own shoes... but obviously, trading makes sense.

    The wealth of nation is based on the division of labour (Adam \noun{Smith}, 1776).

    Typically, we sell just a some products (or services), but we buy a myriad of things and services.

    This is true for individuals, companies, and even countries.
    }
\end{frame}

\begin{frame}{Why Trade?}
    \mytext{
    \begin{itemize}
        \item goods we otherwise would not have
        \begin{itemize}
            \item different climate
            \item different resources
            \item different know-how
        \end{itemize}
        \item goods we have too much of
        \begin{itemize}
            \item dumping
            \item capacities
        \end{itemize}
        \item goods we sell to buy other goods
        \begin{itemize}
            \item (profit) maximisation motive
        \end{itemize}
    \end{itemize}
    }
\end{frame}




\begin{frame}{The Open Economy}
    \showmytikz[1,2]{../shared/tikz/tikz-flowmodel-open}
    \mycomment{
    We understand the economy of a country to consist chiefly of labour markets, the markets for goods and services, and capital markets.

    The state influences in various ways. It taxes its own citizens, regulates activities, provides currency, and purchases from them.
    }
    
\end{frame}

\begin{frame}{The Open Economy}
    \showmytikz[1,2,3]{../shared/tikz/tikz-flowmodel-open}
    \mycomment{
    But the economy is not alone. There are foreign countries, and there is indeed a world market, where there are also markets for labour, goods and services, and capital. 
    
    Because goods can easily be re-routed if there is arbitrage, in effect there is a highly efficient world market. 
    }
    
\end{frame}

\begin{frame}{The Open Economy}
    \showmytikz[2,4]{../shared/tikz/tikz-flowmodel-open}
    \mycomment{
    Unless there are barriers, and abstracting from transportation costs, the markets for goods and services merge. 

    The markets for labour may or may not merge; but labour is clearly less mobile than goods and capital, due to administrative and language barriers, but also due to family bounds.

    There is, however, a connection here too: firms may decide to produce abroad, investing abroad (foreign direct investment, FDI), and (re-)import.

    Capital markets were not always free, but there are no more capital restrictions at present. Monetary flows are also not taxed. The crossing of goods (and services) across borders might be taxed using tariffs or quotas.
    }
\end{frame}

\begin{frame}{Absolute Advantage: I}
    \showmytikz[1]{../shared/tikz/tikz-exchangesmallandsmall}
    \mycomment{
    The continent has more fruit than fish, trading internally at 4:1. So you pay a lot of fruit for one fish.

    The small island has more fish than fruit. Fruit and fish here trade at a rate of 1:4. So on the island, you pay a lot of fish for one fruit.

    Now, if the island were to trade with the continent, every island fisher will be better off by offering their catch on the continent.

    On the continent, they get four fruits for just one fish! That is 16 times as much as on the island! 

    The situation is similar for the continental farmers. 
    }
\end{frame}

\begin{frame}{Absolute Advantage: I}
    \showmytikz[2]{../shared/tikz/tikz-exchangesmallandsmall}
    \mycomment{
    \begin{itemize}
        \item The fruit producer on the island will be outpriced by the continental competition.  The same applies for fishermen on the continent.
        \item In the end, specialisation on both island and continent will be total.
        \item If there are transport costs, these will diminish the effect; but as long as the arbitrage is sufficient, the transport will be paid for.
        \item The island fishermen will compete with the fishermen on the continent; but the fishmongers will work with both, as they are in the same industry. 
        \item Likewise, the former farmers on the island will now probably sell imported fruit.
        \item The traders will ship fish in one direction and fruit in the other.
        \item The new trading rate of fish to fruit is now closer to one, resulting in a healthier diet and larger utility in both industries and both countries.
    \end{itemize}
    }
\end{frame}

\begin{frame}{Absolute Advantage: I}
    \showmytikz[3]{../shared/tikz/tikz-exchangesmallandsmall}
     \mycomment{
    \begin{itemize}
        \item The island farmers will be better off to also start fishing. Island fishing will go up.
        \item There will be more fish on offer on the continent. As a result, the price of fish (relative to fruit) will  fall (and the price of fruit rise, relatively). 
        \item This will have the fishermen on the continent go into the fruit industry, which is now  more lucrative.
        \item In the end, there will be a new equilibrium both on the continent and on the island, with more fish and more fruit on both.  
        \item Fruit producers on the island and fishermen on the continent may (at first) oppose trade. Fishermen on the island and fruit producers on the continent will be in favour.
        \item Both will lose an production industry and become dependent on each other. But both will also have a trading industry, and have both fish and fruit.
    \end{itemize}
    }
\end{frame}

\begin{frame}{Absolute Advantage: II}
    \showmytikz[1]{../shared/tikz/tikz-absoluteadvantage}
    \mycomment{
        Without trade, Red will produce and consume a combination of $X,Y$ at point $A$.

        For example, Red produces more fish ($Y$) and less bread ($X$). That is because for Red it is easy to fish and hard to grow grain.
    }
\end{frame}

\begin{frame}{Absolute Advantage: II}
    \showmytikz[2]{../shared/tikz/tikz-absoluteadvantage}
    \mycomment{
       Without trade, Blue will produce and consume a combination of $X,Y$ at point $B$; so more bread ($X$) than fish ($Y$).
       
       For Blue producing fish ($Y$) is more expensive in absolute terms than producing bread ($X$).  
    }
\end{frame}

\begin{frame}{Absolute  Advantage: II}
    \showmytikz[1,2]{../shared/tikz/tikz-absoluteadvantage}
    \mycomment{

        We see that Red has an absolute advantage in the production of $Y$, and Blue in the production of $X$ relative to one another.
        
        If they trade, each will be happy to offer the respective good at a price the other one cannot hold. 

        Consequentially, they will each focus on whatever it is they can do best.
    }
\end{frame}

\begin{frame}{Absolute Advantage: II}
    \showmytikz[1,2,3]{../shared/tikz/tikz-absoluteadvantage}
    \mycomment{
        If Red and Blue trade, they will enjoy more of $X, Y$  than they could without trade.

        They unite or pool their production frontiers: their coordination results in combined optimisation.

        As a result, each will focus on producing what he/she can do best, that is: cheaper; as he/she can exchange the other good for this. Red will be a fisher, and Blue a farmer.

        In the new equilibrium, both will enjoy a new mix of $X,Y$ that makes them both better off. Red will now eat less fish than before, but more bread, and Blue will enjoy more fish and less bread.

        Total production is likely, but not necessarily, higher than before. But the utility reached under the new distribution will be higher than before. It could also be that both eat better, i.e. a more diverse diet, with less work.
    }
\end{frame}

\begin{frame}{Comparative Advantage: I}
    \showmytikz[1]{../shared/tikz/tikz-exchangelargeandsmall}
    \mycomment{
    Trade and specialisation are a good thing, even if one of the trading countries has an absolute advantage in both industries.

    The rich and large continent has much fish and much fruit; they trade at a rate of 1:1.

    The poor and small island has some fish, but little fruit. Fruit and fish here trade at a rate of 1:4.

    (The numbers here are: 16 fruit and 16 fish on the continent; 1 fruit on the island; 4 fish on the island.)

    Now, if the island were to trade with the continent, every island fisher will still be better off by offering their catch on the continent, because they will get four times as much fruit for it than on the island.
    }
\end{frame}

\begin{frame}{Comparative Advantage: I}
    \showmytikz[2]{../shared/tikz/tikz-exchangelargeandsmall}
    \mycomment{
    \begin{itemize}
        \item The fruit producer on the island will (again) be outpriced by the continental fruit producers. 
        \item Specialisation on the island will (again) be total, but specialisation on the continent only partial.
        \item If there are transport costs, these will diminish the effect; but as long as the arbitrage is sufficient, the transport will be paid for.
        \item The island fishermen will compete with the fishermen on the continent; but the fishmongers will work with both, as they are in the same industry.
        \item The traders will ship both fish and fruit.
        \item The new trading rate of fish to fruit on the island is now also closer to one, resulting in a healthier diet and larger utility for the islanders.
        \item If the island is very small, the impact on the continent is negligible. 
    \end{itemize}
    }
\end{frame}

\begin{frame}{Comparative Advantage: I}
    \showmytikz[3]{../shared/tikz/tikz-exchangelargeandsmall}
     \mycomment{
    \begin{itemize}
        \item The island farmers will be better off to also start fishing. Island fishing will go up.
        \item There will be more fish on offer on the continent (even though not very much more). As a result, the price of fish (relative to fruit) will slightly fall (and the price of fruit rise, relatively). 
        \item This will have some fishermen on the continent go into the fruit industry, which is now (slightly) more lucrative.
        \item In the end, there will be a new equilibrium both on the continent and on the island, with more fish and more fruit on both. The change on the island, however, is more dramatic.
        \item Fruit producers on the island and fishermen on the continent may oppose trade. Fishermen on the island and fruit producers on the continent will be in favour.
        \item Island will lose its fruit industry and become somehow dependent on trade and fruit import.
    \end{itemize}
    }
\end{frame}

\begin{frame}{Comparative Advantage: II}
    \showmytikz[1]{../shared/tikz/tikz-comparativeadvantage}
    \mycomment{
    Now consider the case where blue is just as good at producing $X$ as red, but much worse at producing $Y$.

    Still, trade would be good. If Blue and Red combine their efforts so that Blue does only what they can do best, the trade can enable Red to then do more of what Red can do best.

    The fact that comparative, not absolute, advantage and hence opportuntiy costs of the trading partners alone stipulates trade was first observed by David \noun{Ricardo}.
    }
\end{frame}

\begin{frame}{Comparative Advantage: II}
    \showmytikz[2]{../shared/tikz/tikz-comparativeadvantage}
    \mycomment{
    Now consider the case where blue is just as good at producing $X$ as red, but much worse at producing $Y$.

    Still, trade would be good. If Blue and Red combine their efforts so that Blue does only what they can do best, the trade can enable Red to then do more of what Red can do best.

    The fact that comparative, not absolute, advantage and hence opportunity costs of the trading partners alone stipulates trade was first observed by David \noun{Ricardo}.
    }
\end{frame}

\begin{frame}{Comparative Advantage: II}
    \showmytikz[1,2]{../shared/tikz/tikz-comparativeadvantage}
    \mycomment{
    As we can see, Red has an absolute advantage in both the production of $X$ and $Y$; Blue is nowhere as cheap as Red is. 

    Would you expect that Blue stops producing at all? Of course not. 
    
    When Red and Blue trade, they combine their production frontiers.
    }
\end{frame}

\begin{frame}{Comparative Advantage: II}
    \showmytikz[3]{../shared/tikz/tikz-comparativeadvantage}
    \mycomment{
    We thus arrive at a new, combined production frontier (the new, thick line). 

    Blue will now only do what he can do best and only produce $X$ (at $C$), much of which (but not all) he will trade with Red. 

    Red may only do what he can do best - but maybe he would like to have even more $X$ than Blue can provide, so he may still produce some $X$ (at $D$). So the less productive one will specialise totally, while the more productive one will likely still engage in both activities.

    Note that both will have more of each $X,Y$ than before - and a more balanced diet.

    This shows that with trade, \textbf{everyone should focus on what they can do best.}

    This can lead to the situation where you focus on something (like raw materials) forever.
    
    }
\end{frame}
